\documentclass{article}
\usepackage[preprint]{neurips_2025}
\usepackage[utf8]{inputenc}
\usepackage[T1]{fontenc}
\usepackage{amsmath,amssymb}
\usepackage{microtype}
\usepackage{booktabs}
\usepackage{xurl}
\usepackage{url}
\usepackage[hidelinks]{hyperref}

\title{Reliable Candidate Selection for\\Budget-Constrained Machine Learning Agents}
\author{DATS-SHU-200 Data Mining\\Homework 3: Research Proposal}

\begin{document}
\maketitle

\section{Introduction}

Large language model (LLM) agents can automate parts of machine learning engineering, including writing training code, debugging failures, and improving existing solutions. A central difficulty is deciding which proposed solution deserves an expensive execution. For example, an agent may generate several modifications to a training pipeline, while only a small number can be trained and evaluated within the available time. Executing every proposal consumes resources; discarding proposals prematurely can remove the best solution. This project studies how an agent can use imperfect predictions of candidate quality to make better selection decisions under a limited comparison budget.

Existing agent systems provide two complementary approaches to this problem. AIDE formulates machine learning engineering as tree search over code solutions: the agent drafts, debugs, and improves programs using execution feedback \citep{jiang2025aide}. This makes observed performance available for subsequent decisions, but requires running candidate programs. ForeAgent instead introduces a \emph{Predict-then-Verify} loop \citep{zheng2026predict}. During improvement, it generates multiple candidates and asks an LLM to compare their code using the task description and a data analysis report. A confidence-gated selection procedure then filters candidates before execution. This approach reduces reliance on running every proposal while retaining execution feedback for the candidates that are actually verified.

However, a useful pairwise judgment does not automatically produce a reliable decision among many candidates. Judgments can conflict, and one mistaken comparison can eliminate a strong candidate. Confidence expresses the model's certainty about a particular comparison; it is neither a candidate's performance score nor a guarantee of correctness. The final project asks whether scores from earlier executions in the same run can improve the choice of the next program to train. We will measure the quality of that choice at a fixed comparison budget, then test its effect on the final submission.

The baseline for this project is ForeAgent's candidate-selection component. Its released agent trajectories support a direct check of the original control logic using historical candidate batches and observations. Section~\ref{sec:reproduction} reports three complementary experiments: conditional historical replay, offline selection-quality evaluation, and a single-task end-to-end run. The public Preference Corpus, containing 18,438 candidate pairs across 26 tasks, supplies saved judgments and recorded execution scores for the offline evaluation \citep{zheng2026predict}. MLE-bench provides the broader machine learning engineering task setting \citep{chan2024mle}. Historical-output agreement, prediction accuracy, and end-to-end task performance each supply a different level of evidence.

\section{Baseline Explanation}

\subsection{Problem formulation and notation}

Let $I$ denote a task description, $D$ its dataset, and $C$ a candidate program implementing a machine learning workflow. Let $M(C;I,D)$ denote its task-specific evaluation metric after execution. Define an oriented utility $u(C)$ so that higher values are always better: $u(C)=M(C;I,D)$ for a maximized metric such as AUC, and $u(C)=-M(C;I,D)$ for a minimized metric such as log loss. The underlying engineering objective is to find
\begin{equation}
  C^* \in \operatorname*{arg\,max}_{C\in\mathcal{C}} u(C),
  \label{eq:objective}
\end{equation}
where $\mathcal{C}$ is the space of possible programs. Because obtaining $u(C)$ requires execution, ForeAgent predicts relative quality before spending this cost. This notation describes the optimization objective; the true utility is unavailable to the predictor at decision time.

For candidates $C_i$ and $C_j$, the comparison input consists of $I$, a data analysis report $D_{\mathrm{rep}}$, the two programs, and a prompt $P$. A pretrained LLM with fixed parameters $\theta$ returns
\begin{equation}
  (z_{ij},\widehat{y}_{ij},c_{ij})
  = f_{\theta}(I,D_{\mathrm{rep}},C_i,C_j;P),
  \label{eq:predictor}
\end{equation}
where $z_{ij}$ is the generated reasoning, $\widehat{y}_{ij}\in\{0,1\}$ identifies the preferred input position, and $c_{ij}\in[0,1]$ is the reported confidence. The evaluation label identifies which candidate achieved the better recorded score. The index is relative to the presented pair, so it must be mapped back to candidate identity when comparisons are reused. Crucially, $c_{ij}$ concerns this particular judgment; confidence values from unrelated pairs cannot be treated as a global ranking of candidates.

\subsection{System architecture: prediction followed by verification}

ForeAgent builds on AIDE's search over program nodes. Executed nodes provide code, runtime observations, and measured performance that can inform later drafting, debugging, or improvement. ForeAgent modifies the \emph{improvement} stage by inserting prediction-based filtering between generation and execution \citep[Section~6.2]{zheng2026predict}:
\begin{center}
\small
Executed parent $\longrightarrow$ Candidate generation $\longrightarrow$ LLM comparisons\\[2pt]
$\longrightarrow$ Confidence-gated selection $\longrightarrow$ Execution and feedback.
\end{center}
The published setup proposes $m=10$ improvements to a parent solution, applies a confidence threshold $\tau=0.7$, and selects the top rank group with $k=1$. The selected candidates are subject to the execution policy described below. Actual execution trains and evaluates a candidate program; it is distinct from calling the LLM to compare code. Successful verification returns measured feedback to the search process, so the full agent retains interaction with its computational environment.

The predictor is grounded in task and dataset context rather than code alone. The paper constructs its data reports through a profiling, verification, and verbalization pipeline: generated analysis code is executed, the resulting logs are checked, and the observations are turned into a textual report \citep[Section~3.4]{zheng2026predict}. The report provides information with which the LLM can assess whether a pipeline suits the data. It is part of the prediction input, not the candidate's hidden evaluation score. For the offline baseline, the released reports are reused rather than regenerated.

\subsection{Inference-time selection and execution rules}

The core method uses pretrained LLM inference without fitting an additional reward model or updating model parameters. Accordingly, the relevant technical specification is its decision rule, not a new training loss. The confidence gate accepts a comparison when
\begin{equation}
  a_{ij}=\mathbf{1}\{c_{ij}\geq\tau\},\qquad \tau=0.7.
  \label{eq:gate}
\end{equation}
A high-confidence winner can advance; an inconclusive judgment does not immediately establish a loser. This is a gate on self-reported confidence, rather than a calibrated probability model fitted by this project.

\paragraph{Tournament aggregation.}
In the fixed released implementation, each round compares adjacent active candidates; an odd candidate receives a bye. Confident winners advance and losers are eliminated. Inconclusive pairs enter a pending list. Pending candidates are compared with randomly chosen advanced candidates; if none has advanced, one pending candidate advances directly. A pending winner advances without removing its already-advanced opponent; another inconclusive outcome also advances the pending candidate. The process stops at one survivor or after the round count exceeds $2m$. Final ordering prioritizes survivors, then wins and later rounds; equal consecutive $(\text{wins},\text{round})$ groups share a displayed rank \citep{foreagentcode}.

\paragraph{Selection and probabilistic execution.}
Let $r_i$ be the zero-based displayed rank. The multi-candidate branch marks $r_i<k$ candidates eligible and independently applies an execution probability $p_{\mathrm{exec}}=0.2$ to each. If these tests produce an empty execution set, the original code forces execution of the first rank-0 candidate in sorted order \citep{foreagentcode}. Formally, let $B_i\sim\mathrm{Bernoulli}(p_{\mathrm{exec}})$ and let $i_0$ index that first rank-0 candidate. The preliminary and final execution indicators are
\begin{equation}
  q_i=\mathbf{1}\{r_i<k\}\,B_i,\qquad
  e_i=q_i+\mathbf{1}\!\left\{\sum_jq_j=0\right\}\mathbf{1}\{i=i_0\}.
  \label{eq:execute}
\end{equation}
Thus the original multi-candidate policy executes at least one candidate, while ties may permit several executions. The probability gate and mandatory fallback together specify the execution decision used in Section~\ref{sec:reproduction}.

\subsection{Baseline evaluation and the scope of this project}

The paper evaluates two different capabilities. First, its \emph{Data-centric Solution Preference} task measures how well an LLM predicts pairwise winners. The Preference Corpus contains 895 curated solutions and 18,438 pairs from 26 tasks spanning computer vision, natural language processing, and data science. The principal metric is micro-averaged pairwise accuracy,
\begin{equation}
  \operatorname{Acc}_{\mathrm{pair}}
  =\frac{1}{N}\sum_{\ell=1}^{N}
    \mathbf{1}\{\widehat{y}_{\ell}=y_{\ell}\},
  \label{eq:accuracy}
\end{equation}
where $N$ counts the evaluated pairs and $y_\ell$ is the recorded better-candidate label. The paper reports 61.5\% accuracy for DeepSeek-V3.2-Thinking, averaged over three full-corpus runs \citep[Table~2]{zheng2026predict}. This is a published reference result, not a result achieved by this project.

Second, the full ForeAgent system is evaluated against AIDE on five MLE-bench tasks under a 12-hour limit, with three runs per task. That evaluation measures final task performance through Beat Ratio, the proportion of human leaderboard contestants outperformed, alongside execution-time and search-breadth analyses \citep[Sections~6.3--6.4]{zheng2026predict}. Pairwise accuracy does not establish these end-to-end outcomes: a predictor can make locally correct comparisons yet select poorly within a candidate pool, and a good offline selection does not establish actual training-time savings.

This project treats the pairwise predictor and candidate-selection controller as related baseline components. Earlier exploratory prediction runs used a fixed sample of 246 pairs across 26 tasks. The reproduction reported next evaluates the original ranking and complete execution-control logic at three levels. Historical replay checks agreement with recorded decisions, conditional on archived responses and probability outcomes. Fixed-pool offline evaluation measures the quality of the selected candidates using their recorded scores. A 12-hour Stanford COVID Vaccine run tests the complete generation, selection, execution, and feedback loop. These experiments establish implementation consistency, offline selection quality, and single-run task performance, respectively.

\section{Baseline Reproduction}
\label{sec:reproduction}

\subsection{Objectives and experiment scope}

Three experiments address complementary questions: whether the released controller reproduces archived decisions, how well it selects candidates from fixed pools, and what final task performance it achieves in a live run. All use the original implementation at commit \texttt{c4d52cf} and retain its cache behavior, confidence gate, execution probability, and mandatory execution fallback. The historical experiment is described first; the offline and online results follow in Sections~\ref{sec:offline-selection} and~\ref{sec:online-run}.

The historical reproduction examines whether ForeAgent's released ranking and execution control logic reproduces its recorded historical outputs when supplied with the original candidate batches, generation order, and saved observations. The baseline is the original implementation at commit \texttt{c4d52cf}, with the original cache implementation preserved \citep{foreagentcode}. The comparison backend supplies recorded responses by batch position, and the execution branch receives the probability outcomes recorded in the logs. The resulting experiment is a \emph{conditional historical-trace replay}: it verifies the state updates produced from those observations.

The experiment is frozen as \texttt{foreagent\_original\_historical\_replay\_v1}, dated November~7, 2026. The freeze record binds the source snapshots, input manifests, scripts, verification outputs, and archive using SHA-256 checksums. Subsequent cache corrections and alternative aggregation methods belong to separately versioned experiments. The frozen historical result concerns agreement with the original control logic; the later subsections assess selection quality and live task performance separately.

\subsection{Released records and experimental setup}

The inputs come from the authors' public \texttt{agent\_runs/ForeAgent/} release, fixed at dataset revision \texttt{6b322cb} \citep{foreagentdata}. All 15 released ForeAgent runs were inspected: three runs for each of five MLE-bench tasks. For each run, the inputs include its actual \texttt{config.yaml}, full \texttt{journal.json}, and full \texttt{aide.log}. These 45 original text files were checked against their recorded sizes and hashes. Table~\ref{tab:coverage} summarizes recovery coverage.

\begin{table}[htbp]
  \centering
  \caption{Coverage of the released historical runs. Logged batches include interrupted generation and batches with final-journal gaps. Each recovered batch contains ten candidates.}
  \label{tab:coverage}
  \small
  \begin{tabular}{lrrrr}
    \toprule
    Task & Runs & Logged batches & Recovered & Candidates \\
    \midrule
    Aerial Cactus Identification & 3 & 14 & 12 & 120 \\
    Histopathologic Cancer Detection & 3 & 5 & 4 & 40 \\
    Stanford COVID Vaccine & 3 & 45 & 42 & 420 \\
    Statoil Iceberg Classification & 3 & 60 & 57 & 570 \\
    Ventilator Pressure Prediction & 3 & 16 & 14 & 140 \\
    \midrule
    Total & 15 & 140 & 129 & 1,290 \\
    \bottomrule
  \end{tabular}
\end{table}

The saved configurations agree on the relevant settings: ten improvement candidates, confidence threshold $0.7$, Top-$k=1$, execution probability $0.2$, and a 12-hour agent time limit. The comparison-model label is DeepSeek-V3.2-Thinking with temperature $1.0$. These model settings describe the source runs; this replay consumes their saved responses. Local callbacks record execution decisions, while training and inference observations come from the archived runs. Additional model API calls and candidate-training runs both equal zero.

Recovery uses record availability as the eligibility criterion. Of 140 logged batches, 129 have completed comparison histories and corresponding final-journal nodes. The remaining 11 are retained in a separate coverage ledger: six end during generation, and five contain generated nodes absent from the final journal. This rule was applied using structural completeness, independently of candidate scores or agreement results.

\subsection{Recovery and replay procedure}

\paragraph{Candidate identity, grouping, and order.}
Messages of the form \texttt{Generated improvement candidate i/10} identify the original generation order. Candidate IDs are resolved to their full journal entries and actual source code. The \texttt{node2parent} map verifies a shared parent within each recovered batch. Creation-time ordering provides a second check on the logged sequence. The journal's stored list reflects post-selection saving and re-appending, so generation logs and timestamps establish the replay input order.

\paragraph{Historical comparisons.}
Each candidate in a batch stores the same \texttt{wm\_comparison\_results} sequence, containing winner position, confidence, round, and validity. A local provider supplies these archived results to the original \texttt{tournament\_rank()} calls in their recorded order. The replay checks response counts, round indices, and agreement between confidence and the $0.7$ validity gate. It also checks 8,725 ordinary comparisons against the logged participant IDs and outcomes, and 15,148 revival comparisons against the logged pending-candidate identities and outcomes.

The original controller's random seed or RNG state remains an open recovery item. For revival matches, the advanced opponent is represented by an explicitly marked placeholder choice generated with demonstration seed~0. Responses are supplied by recorded batch position. The checked historical information identifies the pending candidate and outcome; it supplies partial evidence about the advanced opponent. Consequently, the agreement claim concerns control-flow replay conditional on the recorded responses, with exact original revival-pair reconstruction reserved for further log analysis.

\paragraph{Execution decisions and mandatory fallback.}
The original execution statements are extracted directly from \texttt{agent.py}, lines 700--767, and checked for AST equality. Logged probability-trigger outcomes drive the eligible candidates' execution tests. The original code then applies its mandatory fallback: whenever this stage selects an empty execution set, the first rank-0 candidate in sorted order is executed. Local callbacks record selected identities. The resulting execution set is checked against journal \texttt{is\_skipped} flags, and its order is checked against execution-parsing log messages. Supplying observed probability outcomes conditions this check on the historical draws; the verified behavior is Top-$k$ handling, candidate selection, and fallback processing.

\subsection{Results and agreement with the released traces}

For all 129 recovered batches, the original implementation reproduces the recorded candidate ranks, round metadata, sorted order, and final execution decisions. Table~\ref{tab:agreement} reports the comparison units explicitly. Across these batches, the replay consumes 23,873 comparison records, including repeated rounds and cached observations. These records are observations from the historical trajectory, with repeated comparisons contributing to the recorded control flow.

\begin{table}[htbp]
  \centering
  \caption{Agreement conditional on archived comparison responses and logged probability outcomes. Candidate-level and batch-level denominators measure different units.}
  \label{tab:agreement}
  \begin{tabular}{lrr}
    \toprule
    Checked output & Agreement & Rate \\
    \midrule
    Candidate displayed rank & 1,290 / 1,290 & 100\% \\
    Candidate round metadata & 1,290 / 1,290 & 100\% \\
    Complete sorted candidate order & 129 / 129 batches & 100\% \\
    Final execution set & 129 / 129 batches & 100\% \\
    Final execution order & 129 / 129 batches & 100\% \\
    \bottomrule
  \end{tabular}
\end{table}

Mandatory execution fallback occurs in 87 of the recovered batches, matching the log events. For example, in the first recovered batch of the Aerial Cactus run with UUID prefix \texttt{c5d8f2ae}, candidate \texttt{69daaf39} receives rank~0. Its probability test is recorded as skipped, after which the original fallback selects that candidate. Both the resulting execution identity and the final journal flag agree with the replay.

The 100\% values represent \emph{historical-output agreement}. The recorded winner positions are supplied as observations; true candidate-performance scores play no role in recovering ranks. Thus the result establishes faithful processing of these observed histories. Pairwise prediction accuracy, selection quality against executed scores, and end-to-end agent efficiency are separate evaluation questions.

\subsection{Interpretation and reproducibility artifacts}

The replay establishes an implementation baseline for the candidate-selection component: the recovered historical observations lead to the same ranking and execution outputs under the released control rules. The paper's full-corpus prediction accuracy and its five-task online performance results provide broader reference experiments, while this section contributes a direct consistency check against the released agent trajectories \citep{zheng2026predict}. The archived responses already reflect the source runs' inference and cache behavior; fresh model judgments and original revival-opponent reconstruction require their own experimental conditions.

The accompanying \texttt{baseline\_freeze.json} pins the source, data, scripts, and verified outputs by SHA-256. The archive contains full logs and journals, the coverage ledger, and per-batch comparisons. Its \texttt{replay\_author.py} and \texttt{verify\_and\_report.py} scripts reproduce and check these results.


\subsection{Offline selection quality: \texttt{corpus26-v2}}
\label{sec:offline-selection}

\paragraph{Setup.}
This experiment asks whether the original tournament selects candidates with better recorded execution scores than uniform random choice. It uses candidates and one set of published pairwise judgments from the Preference Corpus. Each task contributes 100 candidate-pool instances, each replayed three times with fixed orders and seeds: 2,600 pool instances, 2,164 distinct candidate sets, and 7,800 original-controller replays across 26 tasks. Twenty-one tasks use pools of ten candidates; the other five use pools of 3, 3, 5, 6, and 8. Pool construction uses finite-score availability and pairwise-feedback coverage; score magnitude and candidate rank are reserved for evaluation after selection.

The settings are confidence threshold 0.7, Top-$k=1$, and execution probability 0.2. The original cache is retained and cleared before each replay; comparisons run serially for deterministic replay. Saved judgments replace new LLM responses, and candidate filenames serve as identity stubs in the cache. Actual execution is represented by recording selected identities and looking up their existing scores. Both the probability stage and the original mandatory fallback are included. Additional API calls and candidate-training runs equal zero.

\paragraph{Metrics.}
Let $S$ be the displayed rank-0 group, $B=\operatorname*{arg\,max}_{i\in P}u_i$ the best candidates in pool $P$, and $E$ the final simulated execution set. The selection metric averages a uniform single choice within $S$:
\begin{equation}
  H=\frac{|S\cap B|}{|S|},\qquad
  R=\frac{\max_{i\in P}u_i-|S|^{-1}\sum_{i\in S}u_i}
  {\max_{i\in P}u_i-\min_{i\in P}u_i}.
  \label{eq:offline-metrics}
\end{equation}
Here $R$ is normalized regret for a pool with a nonzero utility range; lower is better. The random reference chooses uniformly from the same pool, with hit probability $|B|/|P|$. We average repeats within a pool, pools within a task, and then the 26 tasks equally. The execution diagnostic is $\mathbf{1}\{E\cap B\ne\varnothing\}$ under the complete original policy. This separates single-choice selection quality from a policy that may execute several tied candidates.

\begin{table}[htbp]
  \centering
  \caption{Original ForeAgent offline selection quality on the fixed corpus pools. Values are task-macro averages. Uniform random choice uses the same candidate pools.}
  \label{tab:offline-quality}
  \begin{tabular}{lrr}
    \toprule
    Metric & Original ForeAgent & Uniform random choice \\
    \midrule
    Uniform-tie best-hit ($H$, higher better) & 20.43\% & 12.53\% \\
    Normalized regret ($R$, lower better) & 0.2810 & 0.3674 \\
    \bottomrule
  \end{tabular}
\end{table}

\paragraph{Results and interpretation.}
Table~\ref{tab:offline-quality} shows a descriptive 7.89 percentage-point advantage in best-hit probability over random choice, together with lower regret. The complete execution simulation includes a best candidate in 21.60\% of replays and executes 1.028 candidates per replay on average. Every replay has a final execution; the mandatory fallback is used in 74.05\% of replays. These values describe this fixed-feedback offline experiment.

The task pools overlap and reuse judgments, so the 7,800 replays are dependent observations. Candidates are assembled from the corpus rather than recovered as siblings generated in one live improvement step. Benchmark judgments and the serial, identity-stub adaptation also differ from fresh judgments under the live Agent prompt. The results therefore support a fixed-pool selection comparison; live search efficiency and robustness across fresh judgments have their own evaluation scope. The original result is retained separately from cache-corrected and alternative-ranking experiments. Evidence is in \texttt{experiments/corpus26-v2/}, including the protocol, complete-execution specification, original results, and verification records. Earlier probability-stage-only execution results are archived as superseded; the table uses the complete-policy version.

\subsection{Single-task end-to-end reproduction}
\label{sec:online-run}

\paragraph{Task and environment.}
Run \texttt{stanford-autodl-v2} executes the original ForeAgent on Stanford COVID Vaccine. The task predicts RNA degradation-related measurements. The prepared data contain 2,160 training sequences and 240 test sequences; the Kaggle archive and four prepared files match the fixed release's recorded MD5 checksums. The released task data report is reused. The runtime uses the pinned author OCI image through PRoot on an AutoDL RTX 4090, with the controller and model requests on a Mac. DeepSeek-V3.2 generates code, DeepSeek-V3.2-Thinking compares candidates, and \texttt{gpt-4o-mini-2024-07-18} parses execution feedback. The DeepSeek requests use the SiliconFlow service with model-name and thinking-mode mapping; the feedback requests use the official OpenAI endpoint.

The run retains ten improvement candidates, confidence 0.70, Top-$k=1$, execution probability 0.20, and the original cache and fallback. The 43,200-second wall-clock budget ran from November~8, 2026, 01:22:02 to 13:22:02, China Standard Time. Generation, API waits and retries, selection, execution, and feedback all consume this budget. The original algorithm selects its final submission by its internal validation metric. External evaluation uses the official private-answer scorer after the run.

\paragraph{Run outputs.}
The completed journal contains 82 nodes: 22 executed and 60 skipped. There are 19 completed outer iterations and seven complete tournament batches. An outer iteration may generate several nodes; these counts use different units. The retained node is \texttt{84c79a8d}, with parsed internal validation value 0.4378. Its submission contains 25,680 valid prediction IDs. The official private-set error is \textbf{0.26640}; an independent, ID-aligned calculation gives 0.2663989312. The scored subset contains 16,320 sequence-position rows. The metric is mean columnwise root mean squared error (MCRMSE) across reactivity, deg\_Mg\_pH10, and deg\_Mg\_50C; lower is better. The internal validation value and private-set score represent separate evaluations.

\begin{table}[htbp]
  \centering
  \caption{Same-task endpoint comparison. The authors' three released final submissions and this run are scored against the same private answers. Each entry is a single run; lower MCRMSE is better.}
  \label{tab:online-score}
  \begin{tabular}{lr}
    \toprule
    Stanford COVID Vaccine submission & Private-set MCRMSE \\
    \midrule
    Author run \texttt{01abb886} & 0.25768 \\
    Author run \texttt{5d12c313} & 0.24320 \\
    Author run \texttt{cbb2a0b0} & 0.27497 \\
    \midrule
    This reproduction (one 12-hour run) & \textbf{0.26640} \\
    \bottomrule
  \end{tabular}
\end{table}

\paragraph{Comparison and operating conditions.}
The score lies within the range of the authors' three same-task endpoints (Table~\ref{tab:online-score}); their mean is 0.25862. The paper evaluates five tasks with three runs per task, whereas this reproduction contributes one run on one of those tasks. This provides evidence that the complete original loop produces a valid submission with a comparable same-task score. Repeated-run variability and the online contribution of filtering require repeated runs and a matched execution-based baseline.

Hardware, model services, and network latency differ from the historical runs. After the run, the experimenter reported that Mac screen sleep and VPN disconnection disrupted model API access. One feedback request, including retries, took 13,491.88 seconds (about 3 hours 45 minutes); this elapsed request time includes processing and retries as well as interruption-related waiting. The entire interval counts toward the standard 12-hour budget. The saved evidence supplies the request elapsed time; the VPN outage duration has a separate measurement gap. This operating condition reduces the time available for productive search and is retained in the reported result. A separately initiated extension was stopped at the user's request with zero new completed nodes; its retained prediction hash matches the standard submission. Table~\ref{tab:online-score} reports the standard 12-hour outcome.

\paragraph{Scoring and artifact checks.}
The original scorer and an independent ID-aligned MCRMSE calculation agree for this submission and all three author submissions. A row-shuffle diagnostic revealed index-sensitive behavior in the original pandas-based scorer: shuffling the same predictions changes its score, while ID-aligned evaluation remains invariant. The archived inputs have been checked for unique IDs, required columns, finite predictions, and agreement between both scoring methods. The original scorer is preserved, and the verification records make its ordering condition explicit.

The source check verifies 484 Git files against the pinned commit. The final prediction file is tied to the retained journal node and archived with its SHA-256 checksum. Run counts, timing, submission hashes, and author references are recorded in \texttt{experiments/online-stanford-v2/}; the independent source, data, artifact, and scoring checks are recorded in \texttt{experiments/environment-acceptance/}. Together with historical replay and fixed-pool evaluation, these results establish three distinct pieces of baseline evidence: faithful control processing, measurable offline selection quality, and a completed single-task Agent run.


\section{Future Plan}
\label{sec:future-memory}

\paragraph{Use earlier executions to choose the next candidate.}
ForeAgent's pairwise judge sees a task description, a data report and two programs. By the time the agent compares a new batch, it may also have trained related programs on this task. We will test whether their validation scores help it choose its next execution. Our control is C0, which retains the original tournament and corrects implementation errors including reversed cache hits. The Stanford result in Section~\ref{sec:online-run} belongs to the original agent.

\paragraph{Method and pilot evidence.}
For each candidate, retrieve the three most similar historical programs by TF--IDF cosine similarity over code tokens and adjacent token pairs. Estimate candidate quality from a similarity-weighted mean of the neighbours' within-history utility ranks. Average the candidate's rank on this estimate with its judge rank. The latter puts the C0 top group first, then orders by accepted win fraction and number of accepted opponents. A pilot used up to ten same-task programs outside each candidate pool and their corpus scores as a proxy for execution history. Across 26 development tasks, fusion reduced task-mean normalized regret from 0.2642 to 0.1966 with DeepSeek judgments and from 0.2725 to 0.2030 with three GPT runs averaged within task. The pilot's primary method, breaking only C0 ties, gained 0.0093 on GPT against a 0.01 advancement criterion. Fusion was its prespecified secondary method.

\paragraph{Priority 1. Check which evidence the agent can use.}
We replayed 136 complete batches from 16 recorded runs. Before generation of each batch began, 130 batches had at least three usable earlier validation scores and 60 had ten. Only six batches had sufficient history and two or more executed siblings with distinct usable outcomes. Their 13 observed pairs span three tasks. Memory alone ordered 29.6\% correctly under task-balanced averaging, compared with 50.8\% after shuffling historical scores. The logged tournament tied on every observed pair. These outcomes describe candidates the original policy actually ran. Evaluating alternative choices across the full pool requires observing programs that policy would skip.

\paragraph{Priority 2. Compare the two choices directly.}
Freeze retrieval and fusion before collecting outcomes. Use up to ten recent valid programs from the same run, with C0 as the fallback when fewer than three are available. Give both selectors the same candidates, comparisons and budget. Record their choices before execution, using a shared seed for ties. When choices differ, train both under a common split, scorer, seed and time cap. Count identical choices as zero paired difference. Report the disagreement rate, within-task score differences, failures and elapsed time. Account for additional evaluation runs separately from each selector's execution cost.

\paragraph{Priority 3. Test the complete agent.}
If fusion improves paired choices across tasks at comparable cost, freeze the method and compare it with C0 on fresh tasks under equal wall-clock limits and paired seeds. Each selector controls its own executions and subsequent search. The final submission score, valid submission rate and total model and training costs will show whether the change improves the complete agent.


\begin{thebibliography}{9}

\bibitem[Jiang et~al.(2025)]{jiang2025aide}
Zhengyao Jiang, Dominik Schmidt, Dhruv Srikanth, Dixing Xu, Ian Kaplan, Deniss Jacenko, and Yuxiang Wu.
\newblock AIDE: AI-Driven Exploration in the Space of Code.
\newblock \emph{arXiv:2502.13138}, 2025.
\newblock \url{https://arxiv.org/abs/2502.13138}.

\bibitem[Zheng et~al.(2026)]{zheng2026predict}
Jingsheng Zheng, Jintian Zhang, Yujie Luo, Yuren Mao, Yunjun Gao, Lun Du, Huajun Chen, and Ningyu Zhang.
\newblock Can We Predict Before Executing Machine Learning Agents?
\newblock In \emph{Proceedings of ACL 2026 (Volume 1: Long Papers)}, pp.\ 3941--3974, 2026.
\newblock \url{https://aclanthology.org/2026.acl-long.182/}.

\bibitem[Chan et~al.(2024)]{chan2024mle}
Jun Shern Chan et~al.
\newblock MLE-bench: Evaluating Machine Learning Agents on Machine Learning Engineering.
\newblock \emph{arXiv:2410.07095}, 2024.
\newblock \url{https://arxiv.org/abs/2410.07095}.

\bibitem[ForeAgent authors(2026a)]{foreagentcode}
ForeAgent authors.
\newblock Released implementation of Predict Before Execute.
\newblock GitHub repository, fixed commit \texttt{c4d52cf} (full identifier in the URL below).
\newblock Relevant files: \texttt{world\_model.py}, \texttt{agent.py}, and \texttt{config.yaml} under \texttt{mle-bench/agents/aide/}.
\newblock \url{https://github.com/zjunlp/predict-before-execute/tree/c4d52cf99bd870d830b456ac7c0684aec1aef375}.


\bibitem[ForeAgent authors(2026b)]{foreagentdata}
ForeAgent authors.
\newblock PredictBeforeExecute: released ForeAgent execution records.
\newblock Hugging Face dataset, fixed revision \texttt{6b322cb}; subdirectory \texttt{agent\_runs/ForeAgent/}.
\newblock \url{https://huggingface.co/datasets/zjunlp/PredictBeforeExecute/tree/6b322cb88bdbcb2b2d3897ec7d0ded94a5bb2d06/agent_runs/ForeAgent}.

\end{thebibliography}

\paragraph{Generative AI use.}
OpenAI Codex assisted with source and log inspection, Python replay and verification tooling, English drafting, and LaTeX typesetting. Requests included (translated from Chinese): ``Replay the original version using the authors' real candidate groups, order, and saved judgments, and check agreement'' and ``Freeze this part and add it to HW3's Baseline Reproduction using the previous NeurIPS template.'' A further request was to incorporate the \texttt{corpus26-v2} offline selection results and the new end-to-end run while retaining the 129-batch historical replay. Original prompts and affected components are recorded in the accompanying \texttt{AI\_USAGE.txt}.

Claude Code assisted with the corpus pilot. OpenAI Codex checked its task-level summaries, implemented the chronological journal audit and drafted Section~\ref{sec:future-memory}. The requests were to read the repository, add the confirmed plan to HW3 for review and begin implementation in priority order. The follow-up instruction was \emph{Choose whichever yields results most easily} (translated from Chinese).
\end{document}
